\section{Introduction}

A graph diagnostic is useful for model selection if following its recommendation improves prediction. Its value depends on the candidate models, the score-to-action rule, and the fallback after abstention. A statistic that favors a graph convolutional network over an MLP may lead to different losses when the graph action selects among six trained architectures~\citep{kipf2017gcn,velivckovic2018gat,zhu2020beyond,lim2021linkx}.

Homophily research has developed class-conditional, feature, structural, and neighborhood-distribution statistics~\citep{luan2023whendo,platonov2023characterizing,zheng2024trihom,pmlr-v235-wang24u}. Learned selectors use graph characteristics and historical performance to choose models on new tasks~\citep{park2023metagl,park2023glemos}. We examine whether a diagnostic's advantage survives a change in candidate portfolio and which datasets determine the aggregate result. We then test whether calibration gains survive holding out related datasets together.

We audit graph-versus-multilayer-perceptron (MLP) decisions on 11 node-classification datasets and ten seeds. Validation selects candidates; oracle-portfolio regret measures the test accuracy lost by the final decision. Benefit and harm relative to always-graph identify useful and costly switches. A frozen nine-policy benchmark establishes the original case: a historical homophily rule with MLP fallback incurs 7.46 percentage points of regret against 0.26 for always-graph. Primary dataset-level tests establish no superiority after multiplicity correction. Separate post-hoc retraining and diagnostic analyses test the boundaries of that result.

A training-label-only Gaussian Naive Bayes adaptation of the Classifier-based Performance Metric (CPM) improves on both constant choices in all 12 restricted GCN/GAT configurations. It loses to always-graph in the four full-portfolio configurations. Roman-empire contributes 78.7--86.5\% of that loss; omitting it reverses three comparisons. Adjusted homophily has lower full-portfolio regret under leave-one-dataset-out (LODO) calibration, but its gains vanish or reverse under source-group holdout. Jointly withholding Cornell and Wisconsin produces an always-graph rule for both datasets in every full-portfolio configuration. Table~\ref{tab:contribution_map} gives the comparisons for one input/budget setting; subsequent sections report all configurations.

\begin{table}[tbp]
\centering
\caption{Connected decision comparisons under CS inputs and 24 MLP trials. Regrets are descriptive percentage points. The graph reference and oracle are recomputed for each portfolio and dataset set. The full 11-dataset set remains primary.}
\label{tab:contribution_map}
\small
\setlength{\tabcolsep}{3pt}
\begin{tabular}{@{}llrr@{}}
\toprule
Decision setting & Selector & Regret & Always-graph \\
\midrule
GCN, all datasets & CPM-.5 & 0.172 & 7.468 \\
Full portfolio, all datasets & CPM-.5 & 1.561 & 0.528 \\
Full, omit Roman-empire & CPM-.5 & 0.232 & 0.580 \\
Full, all datasets & Adjusted homophily, LODO & 0.145 & 0.528 \\
Full, all datasets & Adjusted, source-group holdout & 2.538 & 0.528 \\
\bottomrule
\end{tabular}
\end{table}

The audit extends CPM's paired-classifier evaluation to portfolio decisions and tests adjusted homophily and label informativeness as selection inputs. It complements GLEMOS's learned selection setting by tracing individual switches and their dependence on calibration sources. Released code reconstructs policies from retained candidate records and regenerates bound inputs from public data. The appendices provide mechanistic context through an auxiliary Gaussian calculation and graph intervention.
